\chapter{Agents, Possibilities, and Records}
\label{ch:foundations}

The purpose of the theory is to describe agents whose interactions can produce further agents. The first distinction is between an agent's probability law and the possibilities represented by that law. The second is between comparing two descriptions and specifying an interaction that changes them. A parent construction must preserve enough of both to continue the dynamics.

\section{Statistical sections}

Let $\mathcal C$ be a smooth base of contexts and let $P\to\mathcal C$ be a principal $G$-bundle. A context labels where a description applies; it is not identified with physical position or time. Let $\mathsf K$ and $\mathsf M$ be standard Borel spaces, with measurable $G$-actions by bimeasurable bijections. Their elements $k$ and $m$ represent state contents and model presentations. The actions may be different representations of the same group. They act on probability measures by pushforward.

\definitionheading{Statistical agent section}{def:agent-section}
Set $\mathsf Y=\mathsf K\times\mathsf M$, with the product action, and give $\Prob(\mathsf Y)$ the sigma-algebra generated by evaluation on measurable sets. The associated statistical bundle is
\begin{equation}
\mathcal B=P\times_G\Prob(\mathsf Y),
\qquad (u g,\zeta)\sim(u,g_\#\zeta).
\label{eq:statistical-bundle}
\end{equation}
The statistical component of agent $i$ is a section $\zeta_i$ over a domain $\mathcal U_i\subseteq\mathcal C$. Its full state also includes the internal variables and access mechanisms read by its update rule.

At a fixed context, disintegration gives
\begin{equation}
\zeta_i(dk,dm)=q_i(dk\mid m)s_i(dm),
\qquad q_i(dk)=\int q_i(dk\mid m)s_i(dm).
\label{eq:agent-disintegration}
\end{equation}
Here $q_i$ is the state marginal and $s_i$ is the law over model presentations. Regular conditional laws exist on the stated spaces, with versions defined up to the appropriate marginal null sets \citep{Kallenberg2021}. A model presentation acquires predictive content through a normalized kernel; its name or coordinate value alone supplies no prediction.

The product restriction $\zeta_i=q_i\otimes s_i$ removes dependence between the two sectors within an agent. A product over agents makes a further restriction on represented contents. Neither restriction prevents the update of one law from depending on other laws. Coupling is a property of the dynamics; factorization is a property of a chosen probability description.

Smoothness is an additional hypothesis. When derivatives of laws are required, use a declared regular statistical family stable under the relevant coordinate changes. The whole space $\Prob(\mathsf Y)$ is not assigned a smooth manifold structure by notation. \Cref{ch:geometry} equips an admitted family with the Fisher metric of the full joint law and pulls that metric back to contextual directions through a covariant derivative. The product restriction has a simpler metric, but marginal metrics alone do not measure changes of dependence in a general joint family. A collection of pointwise marginals also does not determine a probability law on fields over $\mathcal C$.

The population at a context consists of the agents whose domains contain that context. Most of the dynamics below is formulated at a common $c_*$. This retains the ambient bundle and local sections without requiring parallel transport between distinct base points. Independent interactions between agents at $c_*$ are defined in \cref{ch:interactions}. Contextual statistical geometry is developed in \cref{ch:geometry}; \cref{app:connection} gives an explicit directional-transport realization of its connection.

\section{Presentation and probability}

In local vector coordinates, write $U_i^q:V_q\to E_i^q$ and $U_i^s:V_s\to E_i^s$ for the sector frames. The spaces $E_i^q,E_i^s$ are the represented state and model fibers. A common principal frame can induce both maps through different representations. Reference laws are
\begin{equation}
\bar q_i=((U_i^q)^{-1})_\#q_i,
\qquad
\bar s_i=((U_i^s)^{-1})_\#s_i.
\label{eq:reference-laws}
\end{equation}
A passive change sends $q_i$ to $(g_i^q)_\#q_i$ and $U_i^q$ to $g_i^qU_i^q$, with the corresponding model-sector rule. The reference laws are unchanged. This is a change of description, not an inference step.

For probability measures $\mu,\nu$ on one measurable space, relative entropy is
\begin{equation}
\KL(\mu\Vert\nu)=\int\log\frac{d\mu}{d\nu}d\mu,
\label{eq:relative-entropy}
\end{equation}
when $\mu\ll\nu$, and is $+\infty$ otherwise. In this formula and below, a displayed integral is understood with its declared measure; no common Lebesgue density is required. Relative entropy compares distinguishability of laws \citep{KullbackLeibler1951}.

\propositionheading{Coordinate invariance}{prop:kl-coordinate}
For a bimeasurable bijection $f$, including divergences with value $+\infty$,
\begin{equation}
\KL(f_\#\mu\Vert f_\#\nu)=\KL(\mu\Vert\nu).
\end{equation}
\begin{proof}
Absolute continuity is preserved in both directions. If $h=d\mu/d\nu$, the transformed density is $h\circ f^{-1}$. Integrating its logarithm against $f_\#\mu$ gives the original integral. If absolute continuity fails, a set of zero $\nu$-measure and positive $\mu$-measure is carried to a set with the same property.
\end{proof}

\definitionheading{Normalized channel}{def:channel}
A Markov kernel $K:X\rightsquigarrow Y$ assigns a probability measure $K(x,\cdot)$ to each $x$, measurably in $x$. Its action on a law is
\begin{equation}
(\mu K)(A)=\int K(x,A)\mu(dx).
\label{eq:kernel-action}
\end{equation}
A deterministic measurable map $f:X\to Y$ is represented by $K_f(x,\cdot)=\delta_{f(x)}$. Then $\mu K_f=f_\#\mu$.

An invertible frame change preserves all distinctions. A selective or noisy channel need not. Both are well-defined operations on general probability laws, but their roles in an agent interaction differ.

\section{Internal records}

At a fixed context, let $\mathsf S_i$ be the admitted internal-state space of agent $i$, and let $\mathsf S$ be the complete configuration space. It contains the statistical laws, source rows, interfaces, and memory variables that the dynamics reads. Write $\operatorname{pr}_i:\mathsf S\to\mathsf S_i$ for the local state projection. For the measurable constructions, take these spaces to be standard Borel. A finite collection of probability-law spaces and finite or Euclidean memory spaces supplies such a realization.

\definitionheading{Record and readout}{def:record}
A record readout is a measurable map
\begin{equation}
\rho_i:\mathsf S_i\longrightarrow\mathsf O_i,
\label{eq:record-readout}
\end{equation}
where $\mathsf O_i$ is the alphabet of internally distinguishable outcomes. If an interaction with source $j$ produces configuration $Z_{t+1}$, its record variable is
\begin{equation}
O_{ij,t}=\rho_i(\operatorname{pr}_i Z_{t+1}).
\label{eq:record-variable}
\end{equation}
For stochastic dynamics this is a measurable function from the probability space of realizations to $\mathsf O_i$. A realized value is denoted $o_{ij,t}$. The source label is part of the record only when the internal state retains it.

A record can be transient. Persistence means that the distinction remains accessible through a specified later readout. A list kept outside the agent does not supply this accessibility. A memory coordinate is therefore an explicit part of the model unless the required distinction is already represented in its state.

\propositionheading{Induced record law}{prop:record-law}
Let $\Gamma_{ij}:\mathsf S\rightsquigarrow\mathsf S$ be a normalized interaction transition. Then
\begin{equation}
\mathcal R_{ij}(x,A)
=\Gamma_{ij}\left(x,\{x':\rho_i(\operatorname{pr}_i x')\in A\}\right),
\label{eq:record-kernel}
\end{equation}
is a normalized record kernel. If $\Gamma_{ij}(x,\cdot)=\delta_{F_{ij}(x)}$, the record is the deterministic function $\rho_i\circ\operatorname{pr}_i\circ F_{ij}$.
\begin{proof}
The right side is the pushforward of a probability measure through the measurable map $\rho_i\circ\operatorname{pr}_i$. Its total mass is one, and measurability in $x$ follows from the defining property of a kernel. The deterministic case follows by evaluating a Dirac measure.
\end{proof}

This definition introduces no observing agent outside the population. It specifies where an outcome is available and how it can enter later dynamics. It does not derive the transition $\Gamma_{ij}$ or decide whether its stochasticity represents objective chance. Those are further modeling choices.

For example, a vector comparison can be represented by
\begin{equation}
r_{ij}:E_i^q\times E_j^q\times\R^r\to\R^r,
\qquad
r_{ij}(k_i,k_j,\epsilon)=A_i k_i-B_{ij}k_j+\epsilon.
\label{eq:sample-comparison}
\end{equation}
A noise law turns this map into a record kernel on represented contents. A joint content law determines its predicted outcomes. A transition on agent states must still specify how an outcome is registered and how the agent changes. A predictive likelihood and an internal update thus have different domains.

\section{What correlations describe}

Suppose two represented bits have the independent uniform law. Register the event that they agree. Their conditional joint law assigns probability $1/2$ to each of $(0,0)$ and $(1,1)$. Each marginal remains uniform, while the conditional expectation of their product changes from $1/4$ to $1/2$. An exact joint posterior has acquired dependence that its two marginals cannot report.

This example identifies a concrete need for a joint law: predicting an observable that reads both bits after the shared record. It does not imply that every admissible agent dynamics must propagate that posterior. A coupled product-family dynamics is another precise model, provided it is declared. Similarly, two routes that inspect the same source must share its identity; replacing it by independent redraws specifies a different experiment.

The configuration $x$ of law-valued agents, a sample tuple $(k_i,m_i)_i$ used in a probabilistic model, and an ensemble of possible configurations are three different objects. A deterministic dynamics on $x$ needs no probability law over configurations. A preparation ensemble can be introduced to assess predictions or compression, but it is then stated explicitly. These distinctions will matter when a parent retains a law over contents while discarding information needed to evolve that law.
